\section{工程规模与综合例题}
\label{l11:sec:examples}

\subsection{两层卷积的形状与运算规模}
\label{l11:subsec:project_scale}

一个用于服饰图像数据集 Fashion MNIST 分类的修改版 LeNet-5 网络，在 Mini-DNN 这一 C++ 神经网络框架中实现，优化对象是使用 CUDA 执行的卷积层前向计算。两次卷积调用分别使用 $86\times86$ 和 $40\times40$ 的输入空间尺寸，批次包含 $100$ 至 $10\,000$ 幅图像，卷积核均为 $7\times7$。这些尺寸指实际传给相应卷积层的张量空间大小。

表\ref{l11:tab:project_shapes}使用步幅 $1$、无填充的定义，列出每幅图像的形状与计算量。参数统计只包含卷积权重，运算统计只包含卷积乘加，均不包含偏置、激活或池化。

\begin{table}[htbp]
\centering
\small
\caption{两次卷积调用的单幅图像规模}
\label{l11:tab:project_shapes}
\begin{tabularx}{\linewidth}{Xrr}
\toprule
项目 & 第一次调用 & 第二次调用\\
\midrule
输入形状 $C\times H\times W$ & $1\times86\times86$ & $4\times40\times40$\\
输出通道数 $M$ & $4$ & $16$\\
核大小 $K\times K$ & $7\times7$ & $7\times7$\\
输出形状 $M\times H_{\mathrm{o}}\times W_{\mathrm{o}}$ & $4\times80\times80$ & $16\times34\times34$\\
卷积权重数 $MCK^2$ & $196$ & $3\,136$\\
输出元素数 $MH_{\mathrm{o}}W_{\mathrm{o}}$ & $25\,600$ & $18\,496$\\
每个输出的乘加数 $CK^2$ & $49$ & $196$\\
每幅图像的乘加数 & $1\,254\,400$ & $3\,625\,216$\\
每幅图像的浮点运算数 & $2\,508\,800$ & $7\,250\,432$\\
\bottomrule
\end{tabularx}
\end{table}

第一次卷积的空间输出为 $86-7+1=80$，第二次为 $40-7+1=34$。如果两次卷积后均采用 $2\times2$、步幅 $2$ 的池化，则空间尺寸分别从 $80$ 降到 $40$、从 $34$ 降到 $17$；这个条件下，第一层池化输出的空间尺寸与第二次卷积输入相容。卷积本身的尺寸和计算量不依赖是否随后采用这一池化配置。

第二次调用的输出元素数仅为第一次的 $18\,496/25\,600=0.7225$，但每个输出需要处理四倍的输入通道数，因此每图计算量比为
\begin{equation}
  \frac{16\cdot34^2\cdot4\cdot7^2}
       {4\cdot80^2\cdot1\cdot7^2}
  =0.7225\times4=2.89.
  \label{l11:eq:project_compute_ratio}
\end{equation}
空间分辨率下降并不保证下一层计算量下降；输入通道数、输出通道数和窗口大小共同决定总量。两次卷积每图合计 $4\,879\,616$ 次乘加，即 $9\,759\,232$ FLOPs。批次为 $B$ 时，这些计算量乘以 $B$，卷积权重数量则保持不变。

\subsection{展开工作区与优化评估}
\label{l11:subsec:project_workspace}

将上述两层改写为矩阵乘法时，第一次的 $D=49,L=6400$，第二次的 $D=196,L=1156$。表\ref{l11:tab:project_workspace}给出显式展开的空间需求，其中单精度浮点格式（FP32）中每个浮点数占 $4$ 字节。

\begin{table}[htbp]
\centering
\small
\caption{显式展开数组的空间需求}
\label{l11:tab:project_workspace}
\begin{tabularx}{\linewidth}{Xrr}
\toprule
项目 & 第一次调用 & 第二次调用\\
\midrule
每幅图像的展开矩阵 $D\times L$ & $49\times6400$ & $196\times1156$\\
每幅图像的展开元素数 & $313\,600$ & $226\,576$\\
每幅图像的 FP32 展开字节数 & $1\,254\,400$ & $906\,304$\\
展开数组与原输入的大小比 & $\approx42.40$ & $35.4025$\\
$B=100$ 的完整展开字节数 & $125\,440\,000$ & $90\,630\,400$\\
$B=10\,000$ 的完整展开字节数 & $12\,544\,000\,000$ & $9\,063\,040\,000$\\
\bottomrule
\end{tabularx}
\end{table}

采用十进制单位 $1\,\mathrm{GB}=10^9$ 字节时，一次展开全部 $10\,000$ 幅图像，两个调用分别需要约 $12.544$ GB 和 $9.06304$ GB 的展开工作区。这些数字仅包含展开数组，不包括输入、权重、输出及其他临时数组。两层按顺序执行并复用工作区时，容量只需覆盖其中较大的需求，不必将二者相加。

完整批次展开是一种执行选择。逐图展开、相乘并复用缓冲区时，工作区可缩小为单幅图像的 $4DL$ 字节；每次处理 $B_c$ 幅图像的小批次时，则为 $4B_cDL$ 字节。小批次降低峰值存储占用，但可能增加执行阶段数量，并改变矩阵乘法的规模与并行度。隐式展开进一步移除了完整展开数组，直接从原输入装入当前计算块。工作区约束因此会影响算法组织，而不仅是最后的内存分配。

实现优化可以按相互依赖的三个技术阶段组织：先建立 CPU 参考结果和 GPU 基础卷积，再通过性能分析定位主要开销并实现矩阵展开及其融合，最后改进卷积与展开中的分块、寻址和存储复用。性能分析（profiling）需要区分实际执行时间、访存开销和资源占用，才能判断某项改动解决了哪一类瓶颈。

比较显式与隐式展开时，应统计完成同一卷积结果所需的全部设备计算阶段。显式方案至少包含展开和矩阵乘法；若另需布局重排，也应计入。只记录矩阵乘法内核的时间，会遗漏显式方案生成输入矩阵的成本。包含或排除主机与设备之间的数据传输也应保持一致。由于内核启动是异步的，计时终点需要对应设备工作实际完成的位置。

正确性核验应覆盖不同批次、多个输入与输出通道、非方形空间尺寸以及不能被线程块尺寸整除的边缘区域。可以先用小整数输入核对展开矩阵和输出的确切值，再用浮点输入比较不同实现的结果。每次改变线程分工、批次布局或求和组织后，都需要重新确认对应的索引关系；性能提升只有在保持预期输出和统计范围一致时才有意义。
